% Fichier engendré par tools/generate-tex.py à partir de 02-fonctions-second-degre.
% Ne pas éditer à la main : tout le texte provient des commentaires du fichier Lean.
\documentclass[11pt,a4paper]{article}

% Compile aussi bien avec pdflatex qu'avec xelatex ou tectonic.
\usepackage{iftex}
\ifPDFTeX
  \usepackage[T1]{fontenc}
  \usepackage[utf8]{inputenc}
\else
  \usepackage{fontspec}
\fi
\usepackage[french]{babel}
\usepackage{amsmath,amssymb,amsthm}
\usepackage[margin=2.5cm]{geometry}
\usepackage[hidelinks]{hyperref}

\theoremstyle{definition}
\newtheorem{definition}{Définition}
\theoremstyle{plain}
\newtheorem{theoreme}{Théorème}
\newtheorem{lemme}[theoreme]{Lemme}

% Renvoi à la déclaration Lean, une ligne après la démonstration, aligné à droite.
\newcommand{\source}[2]{\par\smallskip\noindent\hfill{\footnotesize\href{#1}{\texttt{#2}}}\par\smallskip}

\title{Fonctions, second degré}
\date{}

\begin{document}
\maketitle
% Les identifiants Lean en machine à écrire ne se coupent pas : on tolère des
% espacements plus lâches plutôt que des lignes qui débordent.
\sloppy

\section{FonctionsSecondDegre}

\noindent
Lycée \textemdash{} section \og{} Fonctions, second degré \fg{}.
Énoncés et démonstrations en français : voir FonctionsSecondDegre.tex.

\subsection{Définitions}

\noindent
Toutes les définitions du fichier sont posées ici, avant les démonstrations.

\begin{definition}
Le \emph{trinôme} de coefficients $a$, $b$, $c$ est la fonction
$x \mapsto ax^2 + bx + c$.
\end{definition}
\source{https://github.com/Commutator-IO/learn-lean/blob/main/courses/02-lycee/02-fonctions-second-degre/FonctionsSecondDegre.lean\#L16}{FonctionsSecondDegre.lean\#L16}

\begin{definition}
Le \emph{discriminant} du trinôme est $\Delta = b^2 - 4ac$.
\end{definition}
\source{https://github.com/Commutator-IO/learn-lean/blob/main/courses/02-lycee/02-fonctions-second-degre/FonctionsSecondDegre.lean\#L19}{FonctionsSecondDegre.lean\#L19}

\begin{definition}
La \emph{fonction homographique} est $x \mapsto \dfrac{ax + b}{cx + d}$.
\end{definition}
\source{https://github.com/Commutator-IO/learn-lean/blob/main/courses/02-lycee/02-fonctions-second-degre/FonctionsSecondDegre.lean\#L22}{FonctionsSecondDegre.lean\#L22}

\subsection{Image, antécédent, sens de variation}

\begin{theoreme}
Une fonction croissante conserve l'ordre : si $x \le y$ alors $f(x) \le f(y)$.
\end{theoreme}

\begin{proof}
C'est la définition même de la croissance.
\end{proof}
\source{https://github.com/Commutator-IO/learn-lean/blob/main/courses/02-lycee/02-fonctions-second-degre/FonctionsSecondDegre.lean\#L27}{FonctionsSecondDegre.lean\#L27}

\begin{theoreme}
Composer une fonction croissante avec une fonction décroissante renverse l'ordre.
\end{theoreme}

\begin{proof}
Si $x \le y$, la croissance de $f$ donne $f(x) \le f(y)$, puis la décroissance de $g$
renverse : $g(f(y)) \le g(f(x))$.
\end{proof}
\source{https://github.com/Commutator-IO/learn-lean/blob/main/courses/02-lycee/02-fonctions-second-degre/FonctionsSecondDegre.lean\#L31}{FonctionsSecondDegre.lean\#L31}

\begin{theoreme}
Composer deux fonctions décroissantes donne une fonction croissante.
\end{theoreme}

\begin{proof}
L'ordre est renversé deux fois, donc rétabli.
\end{proof}
\source{https://github.com/Commutator-IO/learn-lean/blob/main/courses/02-lycee/02-fonctions-second-degre/FonctionsSecondDegre.lean\#L35}{FonctionsSecondDegre.lean\#L35}

\subsection{Fonctions de référence}

\begin{theoreme}
La fonction carré est décroissante sur les négatifs et croissante sur les positifs.
\end{theoreme}

\begin{proof}
Pour $x \le y \le 0$, on a $y^2 - x^2 = (y-x)(y+x) \le 0$ car le premier facteur est
positif et le second négatif. Sur les positifs, les deux facteurs sont positifs.
\end{proof}
\source{https://github.com/Commutator-IO/learn-lean/blob/main/courses/02-lycee/02-fonctions-second-degre/FonctionsSecondDegre.lean\#L41}{FonctionsSecondDegre.lean\#L41}

\begin{theoreme}
La fonction inverse est décroissante sur les réels strictement positifs.
\end{theoreme}

\begin{proof}
Pour $0 < x \le y$, comparer $1/x$ et $1/y$ revient, après multiplication par $xy > 0$,
à comparer $y$ et $x$.
\end{proof}
\source{https://github.com/Commutator-IO/learn-lean/blob/main/courses/02-lycee/02-fonctions-second-degre/FonctionsSecondDegre.lean\#L52}{FonctionsSecondDegre.lean\#L52}

\begin{theoreme}
La fonction racine carrée est croissante.
\end{theoreme}

\begin{proof}
La racine carrée conserve l'ordre sur les positifs, deux nombres positifs étant rangés
comme leurs carrés.
\end{proof}
\source{https://github.com/Commutator-IO/learn-lean/blob/main/courses/02-lycee/02-fonctions-second-degre/FonctionsSecondDegre.lean\#L59}{FonctionsSecondDegre.lean\#L59}

\begin{theoreme}
La fonction cube est strictement croissante sur tout $\mathbb{R}$, contrairement à la
fonction carré.
\end{theoreme}

\begin{proof}
Une puissance d'exposant impair est strictement croissante sur $\mathbb{R}$ tout entier :
elle conserve le signe, là où le carré l'efface.
\end{proof}
\source{https://github.com/Commutator-IO/learn-lean/blob/main/courses/02-lycee/02-fonctions-second-degre/FonctionsSecondDegre.lean\#L63}{FonctionsSecondDegre.lean\#L63}

\subsection{Forme canonique}

\begin{theoreme}
Forme canonique : pour $a \ne 0$,
\[ ax^2 + bx + c = a\left(x + \frac{b}{2a}\right)^2 - \frac{\Delta}{4a} . \]
\end{theoreme}

\begin{proof}
On met $a$ en facteur puis on complète le carré :
\[ ax^2 + bx + c = a\left(x^2 + \frac{b}{a}x\right) + c
   = a\left(x + \frac{b}{2a}\right)^2 - \frac{b^2}{4a} + c , \]
et $-\dfrac{b^2}{4a} + c = -\dfrac{b^2 - 4ac}{4a} = -\dfrac{\Delta}{4a}$.
\end{proof}
\source{https://github.com/Commutator-IO/learn-lean/blob/main/courses/02-lycee/02-fonctions-second-degre/FonctionsSecondDegre.lean\#L69}{FonctionsSecondDegre.lean\#L69}

\subsection{Discriminant et racines}

\begin{theoreme}
Si $\Delta > 0$, le trinôme a deux racines distinctes,
\[ x_{1,2} = \frac{-b \pm \sqrt{\Delta}}{2a} . \]
\end{theoreme}

\begin{proof}
Posons $d = \sqrt{\Delta} > 0$, de sorte que $d^2 = \Delta$. Les deux nombres proposés
sont distincts puisque $d \ne 0$, et en substituant dans le trinôme, le calcul se réduit à
$d^2 - \Delta = 0$.
\end{proof}
\source{https://github.com/Commutator-IO/learn-lean/blob/main/courses/02-lycee/02-fonctions-second-degre/FonctionsSecondDegre.lean\#L78}{FonctionsSecondDegre.lean\#L78}

\begin{theoreme}
Si $\Delta < 0$, le trinôme ne s'annule jamais : il garde le signe de $a$.
\end{theoreme}

\begin{proof}
La forme canonique écrit le trinôme comme $a\,(\cdot)^2 - \Delta/(4a)$. Pour $a > 0$, le
premier terme est positif et $-\Delta/(4a)$ est strictement positif : la somme ne peut pas
être nulle. Pour $a < 0$, les deux termes sont négatifs, même conclusion.
\end{proof}
\source{https://github.com/Commutator-IO/learn-lean/blob/main/courses/02-lycee/02-fonctions-second-degre/FonctionsSecondDegre.lean\#L96}{FonctionsSecondDegre.lean\#L96}

\begin{theoreme}
Si $\Delta = 0$, le trinôme a une racine double en $-\dfrac{b}{2a}$, et s'écrit
$a\left(x + \dfrac{b}{2a}\right)^2$.
\end{theoreme}

\begin{proof}
La forme canonique perd son second terme, et le carré s'annule exactement en
$-b/(2a)$.
\end{proof}
\source{https://github.com/Commutator-IO/learn-lean/blob/main/courses/02-lycee/02-fonctions-second-degre/FonctionsSecondDegre.lean\#L111}{FonctionsSecondDegre.lean\#L111}

\subsection{Factorisation et signe du trinôme}

\begin{theoreme}
Factorisation par les racines : si $x_1 + x_2 = -b/a$ et $x_1x_2 = c/a$, alors
\[ ax^2 + bx + c = a(x - x_1)(x - x_2) . \]
\end{theoreme}

\begin{proof}
Les deux hypothèses donnent $b = -a(x_1 + x_2)$ et $c = a\,x_1x_2$ ; en les reportant et
en développant le membre de droite, les deux expressions coïncident.
\end{proof}
\source{https://github.com/Commutator-IO/learn-lean/blob/main/courses/02-lycee/02-fonctions-second-degre/FonctionsSecondDegre.lean\#L125}{FonctionsSecondDegre.lean\#L125}

\begin{theoreme}
Signe du trinôme : pour $a > 0$ et $x_1 < x_2$, il est négatif entre les racines et
positif à l'extérieur.
\end{theoreme}

\begin{proof}
Entre les racines, $x - x_1 > 0$ et $x - x_2 < 0$ : le produit est négatif, et
multiplier par $a > 0$ ne change pas le signe. À l'extérieur, les deux facteurs sont de
même signe, donc leur produit est positif. D'où la règle : « du signe de $a$ sauf entre les
racines ».
\end{proof}
\source{https://github.com/Commutator-IO/learn-lean/blob/main/courses/02-lycee/02-fonctions-second-degre/FonctionsSecondDegre.lean\#L140}{FonctionsSecondDegre.lean\#L140}

\subsection{Somme et produit des racines}

\begin{theoreme}
Somme et produit des racines : si le trinôme se factorise en $a(x-x_1)(x-x_2)$, alors
\[ x_1 + x_2 = -\frac{b}{a}, \qquad x_1 x_2 = \frac{c}{a} . \]
\end{theoreme}

\begin{proof}
On évalue l'identité en trois points. En $0$ elle donne $c = a\,x_1x_2$, d'où le
produit. En $1$ et en $-1$, la différence des deux égalités élimine $a$ et $c$ et laisse
$2b = -2a(x_1 + x_2)$, d'où la somme.
\end{proof}
\source{https://github.com/Commutator-IO/learn-lean/blob/main/courses/02-lycee/02-fonctions-second-degre/FonctionsSecondDegre.lean\#L162}{FonctionsSecondDegre.lean\#L162}

\subsection{Sommet de la parabole et axe de symétrie}

\begin{theoreme}
Le sommet de la parabole est atteint en $-\dfrac{b}{2a}$, et c'est un minimum lorsque
$a > 0$.
\end{theoreme}

\begin{proof}
Dans la forme canonique, le carré s'annule en $-b/(2a)$ et est positif ailleurs ;
multiplié par $a > 0$, il ne peut qu'augmenter la valeur.
\end{proof}
\source{https://github.com/Commutator-IO/learn-lean/blob/main/courses/02-lycee/02-fonctions-second-degre/FonctionsSecondDegre.lean\#L179}{FonctionsSecondDegre.lean\#L179}

\begin{theoreme}
La parabole est symétrique par rapport à la droite verticale passant par le sommet :
\[ f\left(-\frac{b}{2a} + t\right) = f\left(-\frac{b}{2a} - t\right) . \]
\end{theoreme}

\begin{proof}
Dans la forme canonique, les deux abscisses donnent le même carré $t^2$, donc la même
image.
\end{proof}
\source{https://github.com/Commutator-IO/learn-lean/blob/main/courses/02-lycee/02-fonctions-second-degre/FonctionsSecondDegre.lean\#L189}{FonctionsSecondDegre.lean\#L189}

\subsection{Fonction homographique}

\begin{theoreme}
La fonction homographique n'est pas définie là où son dénominateur s'annule, c'est-à-dire
en $x = -d/c$.
\end{theoreme}

\begin{proof}
En cette valeur, $cx + d = 0$ : le quotient n'a pas de sens. C'est la valeur interdite,
qui donne l'asymptote verticale.
\end{proof}
\source{https://github.com/Commutator-IO/learn-lean/blob/main/courses/02-lycee/02-fonctions-second-degre/FonctionsSecondDegre.lean\#L199}{FonctionsSecondDegre.lean\#L199}

\begin{theoreme}
Elle s'écrit comme la somme d'une constante et d'un terme qui tend vers zéro :
\[ \frac{ax+b}{cx+d} = \frac{a}{c} + \frac{bc - ad}{c\,(cx+d)} , \]
ce qui donne l'asymptote horizontale $y = a/c$.
\end{theoreme}

\begin{proof}
On réduit le membre de droite au même dénominateur $c(cx+d)$ :
\[ \frac{a(cx+d) + bc - ad}{c(cx+d)} = \frac{acx + bc}{c(cx+d)} = \frac{ax+b}{cx+d} . \]
Le second terme tend vers zéro quand $x$ grandit, d'où l'asymptote.
\end{proof}
\source{https://github.com/Commutator-IO/learn-lean/blob/main/courses/02-lycee/02-fonctions-second-degre/FonctionsSecondDegre.lean\#L206}{FonctionsSecondDegre.lean\#L206}

\subsection{Parité}

\begin{theoreme}
Une fonction paire a une courbe symétrique par rapport à l'axe des ordonnées.
\end{theoreme}

\begin{proof}
C'est la traduction de $f(-x) = f(x)$ : deux abscisses opposées ont la même
ordonnée.
\end{proof}
\source{https://github.com/Commutator-IO/learn-lean/blob/main/courses/02-lycee/02-fonctions-second-degre/FonctionsSecondDegre.lean\#L215}{FonctionsSecondDegre.lean\#L215}

\begin{theoreme}
Une fonction impaire a une courbe symétrique par rapport à l'origine.
\end{theoreme}

\begin{proof}
C'est la traduction de $f(-x) = -f(x)$ : deux abscisses opposées ont des ordonnées
opposées.
\end{proof}
\source{https://github.com/Commutator-IO/learn-lean/blob/main/courses/02-lycee/02-fonctions-second-degre/FonctionsSecondDegre.lean\#L219}{FonctionsSecondDegre.lean\#L219}

\begin{theoreme}
La fonction carré est paire, la fonction cube est impaire.
\end{theoreme}

\begin{proof}
$(-x)^2 = x^2$ et $(-x)^3 = -x^3$, l'exposant pair effaçant le signe, l'exposant impair
le conservant.
\end{proof}
\source{https://github.com/Commutator-IO/learn-lean/blob/main/courses/02-lycee/02-fonctions-second-degre/FonctionsSecondDegre.lean\#L223}{FonctionsSecondDegre.lean\#L223}

\vfill
\noindent\rule{0.35\linewidth}{0.4pt}\par
\noindent{\footnotesize Leonardo de Moura et Sebastian Ullrich,
\emph{The Lean~4 Theorem Prover and Programming Language},
\emph{Automated Deduction — CADE~28}, Springer, 2021, p.~625--635,
\href{https://doi.org/10.1007/978-3-030-79876-5_37}{doi:10.1007/978-3-030-79876-5\_37}.\par}

\end{document}
